\scanpage{122}

\[
\forall(s,t)\in S,\ \text{on a}\qquad
 \{s,t\}\in A_0\Longleftrightarrow\{s,\lambda t\}\in A_1.
\tag*{$(45)$}
\]

On voit que cela implique que \deleted{\unsure{la}} les groupes d'autom. $G_0,G_1$ de $P_0,P_1$ sont égaux, et que par suite $\lambda$ \inserted{donc $\lambda^{-1}$} commute aussi à $G_{1}$, et que \deleted{\unsure{l'on|l'ensemble}} par suite $A_0$ est ``transformé'' de $A_1$ par l'élément $\lambda^{-1}$ du commutant de $G_1$ dans $\unsure{\mathfrak S_S|G_S}$ : la relation envisagée est donc symétrique, elle est évidemment réflexive, et on voit tout de suite qu'elle est également transitive : \deleted{elle} c'est donc bien une \uline{relation d'équivalence} entre structures polyédrales spéciales. Si \deleted{\unsure{$C$}} \inserted{$\Pi$} est une classe d'équivalence de telles structures, il lui correspond un groupe

% slanted
\margin{on note l'ens. $\Pi$ \unsure{indiciellement|radicalement}}

\[
\begin{aligned}
 G_0,\quad &\text{groupe d'automorphismes}\\
 &\text{commun aux structures}\\
 &\deleted{\unsure{F}}P_i\quad\deleted{\unsure{\in C}}\quad(i\in\Pi)
\end{aligned}\tag*{$(46)$}
\]

et un groupe

\[
 Z'=\operatorname{Comm}_{\mathfrak S_S}(G_0)\ ,
\tag*{$(47)$}
\]

enfin, on voit que le groupe $Z'$ opère sur $\Pi$, transitivement par définition
% archive sheet number: 121
